%% ============================================================
%% Embedded Module Architecture 実装ガイド
%% ============================================================
%% ビルド: latexmk main.tex (.latexmkrc で uplatex + dvipdfmx)
%% ============================================================
\documentclass[a4paper,11pt,dvipdfmx]{jsarticle}

%% ─── パッケージ ───────────────────────────────
\usepackage[dvipdfmx]{graphicx}
\usepackage[dvipdfmx]{xcolor}
\usepackage[dvipdfmx,
  bookmarks=true,
  bookmarksnumbered=true,
  colorlinks=true,
  linkcolor={blue!70!black},
  citecolor={green!50!black},
  urlcolor={blue!60!black},
  pdfauthor={Shiozawa},
  pdftitle={Embedded Module Architecture 実装ガイド},
  pdfsubject={モジュール設計パターンの実装例と解説},
]{hyperref}
\usepackage{pxjahyper}
\usepackage[top=25mm,bottom=25mm,left=20mm,right=20mm]{geometry}
\usepackage{fancyhdr}
\usepackage{lastpage}
\usepackage{tabularx}
\usepackage{booktabs}
\usepackage{longtable}
\usepackage{enumitem}
\usepackage{titlesec}
\usepackage{tcolorbox}
\tcbuselibrary{skins,breakable}
\tcbset{before={\par\medskip\needspace{4\baselineskip}\noindent}}  % ボックス前に最低4行分確保（タイトル孤立防止）
\usepackage{tikz}
\usepackage{listings}
\usepackage{plistings}
\usepackage{caption}
\usepackage{float}
\usepackage{needspace}
\usepackage{placeins}
\usepackage{amssymb}            % checkmarkシンボル

%% ─── 改ページ制御 ─────────────────────────────
% 孤立行（widow/orphan）防止 — ページ末尾/先頭に1行だけ残さない
\widowpenalty=10000
\clubpenalty=10000
% 数式前後の改ページ抑制
\predisplaypenalty=10000
% 脚注の分割禁止
\interfootnotelinepenalty=10000
% 改ページコスト調整 — 見出し直後や短い段落での改ページを避ける
\brokenpenalty=5000

%% ─── 色定義 ──────────────────────────────────
\definecolor{headerblue}{HTML}{1B3A5C}
\definecolor{accentblue}{HTML}{2E86C1}
\definecolor{lightgray}{HTML}{F2F3F4}
\definecolor{codebg}{HTML}{F7F9FB}
\definecolor{warnred}{HTML}{E74C3C}
\definecolor{noteyellow}{HTML}{F39C12}
\definecolor{okgreen}{HTML}{27AE60}

%% ─── ヘッダー・フッター ───────────────────────
\pagestyle{fancy}
\fancyhf{}
\fancyhead[L]{\small\color{headerblue}\textbf{Embedded Module Architecture} 実装ガイド}
\fancyhead[R]{\small\color{headerblue}Rev.\,1.1}
\fancyfoot[C]{\small\thepage\,/\,\pageref{LastPage}}
\fancyfoot[R]{\small\color{gray}2026-03-17}
\renewcommand{\headrulewidth}{0.8pt}
\renewcommand{\footrulewidth}{0.4pt}

%% ─── セクション見出しスタイル ─────────────────
\titleformat{\section}
  {\needspace{5\baselineskip}\Large\bfseries\color{headerblue}}
  {\thesection}{1em}{}
  [\vspace{-0.5em}\textcolor{accentblue}{\rule{\textwidth}{1.2pt}}]
\titleformat{\subsection}
  {\needspace{4\baselineskip}\large\bfseries\color{headerblue}}
  {\thesubsection}{1em}{}
\titleformat{\subsubsection}
  {\needspace{3\baselineskip}\normalsize\bfseries\color{accentblue}}
  {\thesubsubsection}{1em}{}
\titleformat{\paragraph}[runin]
  {\normalsize\bfseries\color{headerblue}}
  {}{}{}[　]

%% ─── listings 設定 (C++) ──────────────────────
\lstdefinestyle{cppstyle}{
  language=C++,
  basicstyle=\ttfamily\small,
  keywordstyle=\color{accentblue}\bfseries,
  commentstyle=\color{gray}\itshape,
  stringstyle=\color{okgreen},
  numberstyle=\tiny\color{gray},
  numbers=left,
  numbersep=8pt,
  frame=single,
  framerule=0.5pt,
  rulecolor=\color{accentblue!40},
  backgroundcolor=\color{codebg},
  breaklines=true,
  breakatwhitespace=false,
  tabsize=4,
  showstringspaces=false,
  captionpos=b,
  xleftmargin=15pt,
  framexleftmargin=15pt,
  morekeywords={override,nullptr,uint8_t,uint16_t,uint32_t,int32_t,int8_t,bool,volatile,size_t},
}
\lstset{style=cppstyle}

%% ─── tcolorbox スタイル ──────────────────────
% パターン解説ボックス
\newtcolorbox{patternbox}[1]{
  enhanced,
  colback=white,
  colframe=accentblue,
  coltitle=white,
  fonttitle=\bfseries\large,
  title=#1,
  top=4mm,
  arc=2mm,
  breakable,
  before upper={\parindent0pt},
  attach boxed title to top left={yshift=-3mm,xshift=5mm},
  boxed title style={colback=accentblue,arc=1.5mm},
}

% ポイント解説ボックス
\newtcolorbox{pointbox}[1][ポイント]{
  enhanced,
  colback=okgreen!5,
  colframe=okgreen!60,
  coltitle=white,
  fonttitle=\bfseries,
  title={$\checkmark$ #1},
  arc=1.5mm,
  breakable,
  attach boxed title to top left={yshift=-2mm,xshift=5mm},
  boxed title style={colback=okgreen!60,arc=1mm},
}

% 注意ボックス
\newtcolorbox{cautionbox}[1][注意]{
  enhanced,
  colback=warnred!5,
  colframe=warnred!70,
  coltitle=white,
  fonttitle=\bfseries,
  title={! #1},
  arc=1.5mm,
  breakable,
  attach boxed title to top left={yshift=-2mm,xshift=5mm},
  boxed title style={colback=warnred!70,arc=1mm},
}

% ノートボックス
\newtcolorbox{notebox}[1][補足]{
  enhanced,
  colback=accentblue!5,
  colframe=accentblue!50,
  coltitle=white,
  fonttitle=\bfseries,
  title={$\triangleright$ #1},
  arc=1.5mm,
  breakable,
  attach boxed title to top left={yshift=-2mm,xshift=5mm},
  boxed title style={colback=accentblue!50,arc=1mm},
}

% 機能ブロック
\newtcolorbox{funcblock}[1]{
  enhanced,
  colback=white,
  colframe=accentblue!25,
  coltitle=headerblue,
  fonttitle=\bfseries\ttfamily\small,
  title={◆ #1},
  arc=1mm,
  boxrule=0.4pt,
  breakable,
  left=3mm, right=3mm, top=2mm, bottom=2mm,
  before skip=3mm, after skip=3mm,
  before upper={\parindent0pt},
}

%% ─── ドキュメント情報 ─────────────────────────
\newcommand{\doctitle}{Embedded Module Architecture\\実装ガイド}
\newcommand{\docsubtitle}{リポジトリの各モジュールで学ぶ設計パターンの実践}
\newcommand{\docauthor}{Shiozawa}
\newcommand{\docversion}{1.1}
\newcommand{\docdate}{2026年3月17日}

%% ============================================================
\begin{document}

% 表紙
\begin{titlepage}
\begin{center}
\vspace*{30mm}

{\color{accentblue}\rule{\textwidth}{2pt}}
\vspace{8mm}

{\Huge\bfseries\color{headerblue} \doctitle}

\vspace{6mm}

{\Large\color{accentblue} \docsubtitle}

\vspace{4mm}
{\color{accentblue}\rule{\textwidth}{2pt}}

\vspace{20mm}

\begin{tabular}{ll}
    \textbf{バージョン} & \docversion \\[4pt]
    \textbf{作成日} & \docdate \\[4pt]
    \textbf{著者} & \docauthor \\[4pt]
\end{tabular}

\vspace{30mm}

\begin{tcolorbox}[
  enhanced,
  colback=accentblue!5,
  colframe=accentblue,
  arc=2mm,
  width=14cm,
  breakable,
]
    \centering
    本ドキュメントは、Embedded Module Architecture リポジトリに\\
    実装されている各モジュールを題材に、\\
    設計パターンの適用方法を具体的なコードで解説するガイドです。\\[6pt]
    設計仕様書（設計ルール）・教科書（パターンの学習）と併せてお読みください。
\end{tcolorbox}

\vfill
\end{center}
\end{titlepage}

% 改訂履歴
\section*{改訂履歴}
\addcontentsline{toc}{section}{改訂履歴}

\begin{tabularx}{\textwidth}{c c c X}
\toprule
\textbf{版} & \textbf{日付} & \textbf{著者} & \textbf{変更内容} \\
\midrule
1.0 & 2026-03-15 & Shiozawa & 初版作成。全14モジュールの実装解説 \\
1.1 & 2026-03-17 & Shiozawa & 関連ドキュメントの相互参照を追加。日付・Rev番号を全ドキュメントと同期 \\
\bottomrule
\end{tabularx}

\newpage
\tableofcontents
\newpage
% 概要
\section{この資料について}

\subsection{位置づけ}

本リポジトリには3つのドキュメントがある。

\begin{tabularx}{\textwidth}{l l X}
\toprule
\textbf{ドキュメント} & \textbf{ディレクトリ} & \textbf{内容} \\
\midrule
設計仕様書 & \texttt{doc/} & アーキテクチャのルール・API仕様（「何を守るか」） \\
教科書 & \texttt{doc-textbook/} & パターンの段階的な学習（「なぜそうするか」） \\
\textbf{実装ガイド}（本書） & \texttt{doc-guide/} & 実コードの解説（「実際にどう書いたか」） \\
\bottomrule
\end{tabularx}

\begin{cautionbox}[サンプルコードについて]
本リポジトリに実装されているモジュール群のコードは、あくまで\textbf{アーキテクチャの設計パターンを具体的に示すためのサンプル実装}である。
実機での動作検証は行っていないため、実プロジェクトへの適用時は十分なテストを実施すること。
\end{cautionbox}

\subsection{読み方}

本書では、リポジトリに実装されている各モジュールを取り上げ、以下の構成で解説する。

\begin{enumerate}
    \item \textbf{概要} — モジュールの役割と、そのモジュールで学べるパターン
    \item \textbf{コード解説} — 実際のコードを引用し、設計ルールとの対応を示す
    \item \textbf{ポイント} — 特に注目すべき設計判断やテクニック
\end{enumerate}

モジュールは簡単なものから順に並べている。
前半のモジュールで基本パターンを押さえたうえで、後半の応用パターンに進む構成になっている。

\subsection{モジュール一覧}

\begin{longtable}{l l l}
\toprule
\textbf{モジュール} & \textbf{種別} & \textbf{主な学習パターン} \\
\midrule
\endhead
LedModule & 出力 & 最小構成の基本パターン \\
ButtonModule & 入力 & デバウンス（ModuleTimer活用） \\
BatteryModule & 入力 & ADC読み取り + 移動平均フィルタ \\
ServoModule & 出力 & PWM制御（LEDC） \\
BuzzerModule & 出力 & deinit()によるリソース解放 \\
Mpu6500Module & 入力 & I2Cバスパターン + レジスタ直接アクセス \\
TouchModule & 入力 & バスオブジェクト共有（TFT\_eSPIドライバ） \\
TftModule & 出力 & SPI共有 + 複合デバイスの分離 \\
CameraModule & 入力 & PSRAMバッファ + フレーム管理 \\
DriveMotorModule & 出力 & PWM 2ピン制御 + 統合モジュール部品 \\
ChassisModule & 出力 & 機能統合系（内部にモジュールを持つ） \\
BleModule & 入力 & スレッドセーフティ（volatileフラグ方式） \\
WifiModule & 入力 & ステートマシン内蔵 \\
EncoderModule & 入力 & ハードウェア割り込み連携 \\
\bottomrule
\end{longtable}

% LedModule
\section{【基本】LedModule — 最小構成の基本パターン}

\begin{patternbox}{LedModule で学べるパターン}
\begin{itemize}
    \item Config構造体によるピンアサインの外部注入
    \item Data構造体のデフォルト値
    \item .h/.cpp分離（循環依存回避）
    \item 出力モジュールの基本形（SystemDataの値に従って動作するだけ）
\end{itemize}
\end{patternbox}

\subsection{ヘッダー（LedModule.h）}

\begin{lstlisting}[style=cppstyle, caption={LedModule.h}]
#pragma once
#include <Arduino.h>
#include "IModule.h"

struct LedConfig {
    uint8_t ledPin;  // LED出力ピン
};

struct LedData {
    bool state = false;       // LED状態（デフォルト: 消灯）
    uint8_t brightness = 0;   // 輝度（将来拡張用）
};

class LedModule : public IModule {
private:
    LedConfig _config;
public:
    LedModule(const LedConfig& config);
    bool init() override;
    void update(SystemData& data) override;
};
\end{lstlisting}

\begin{pointbox}[Config構造体のポイント]
\texttt{LedConfig}にはピン番号だけを持つ。
ピン番号を\texttt{LedModule}内にハードコードせず、Config構造体で外部から注入する。
これにより、同じ\texttt{LedModule}クラスを別プロジェクトでピンだけ変えて再利用できる。
\end{pointbox}

\begin{pointbox}[Data構造体のデフォルト値]
\texttt{LedData}の\texttt{state = false}は「\texttt{update()}が1度も呼ばれていない状態でも安全」を保証する。
ロジックフェーズが\texttt{state}を\texttt{true}に変更するまで、LEDは消灯したままとなる。
\end{pointbox}

\subsection{実装（LedModule.cpp）}

\begin{lstlisting}[style=cppstyle, caption={LedModule.cpp}]
#include "LedModule.h"
#include "SystemData.h"  // ヘッダーでは前方宣言のみ、cppで完全な定義を取得

LedModule::LedModule(const LedConfig& config) : _config(config) {}

bool LedModule::init() {
    pinMode(_config.ledPin, OUTPUT);
    Serial.println("[Led] init OK");
    return true;
}

void LedModule::update(SystemData& data) {
    digitalWrite(_config.ledPin, data.led.state ? HIGH : LOW);
}
\end{lstlisting}

\begin{pointbox}[出力モジュールの責務]
\texttt{update()}は\textbf{自身のData構造体の値に従って動作するだけ}。
「いつ点灯するか」「何の条件で消すか」などのロジックは一切書かない。
それはロジックフェーズ（\texttt{applyPattern()}等）の仕事である。
\end{pointbox}

\begin{pointbox}[.h/.cpp分離のルール]
ヘッダーでは\texttt{SystemData.h}をインクルードせず、\texttt{IModule.h}内の前方宣言だけを利用する。
\texttt{SystemData.h}のインクルードは\texttt{.cpp}側で行う。
これにより\texttt{LedModule.h} → \texttt{SystemData} → \texttt{LedModule.h}の循環依存を防ぐ。
\end{pointbox}

\subsection{ProjectConfig.hでの定義}

\begin{lstlisting}[style=cppstyle, caption={ProjectConfig.hでのConfig定義}]
const LedConfig LED_CONFIG = {
    .ledPin = 2,  // GPIO2: オンボードLED
};
\end{lstlisting}

回路設計者はこのファイルだけ見ればピンアサインが分かる。
\texttt{LED\_CONFIG}という大文字スネークケースの命名規則も守っている。

% ButtonModule
\section{【基本】ButtonModule — デバウンスパターン}

\begin{patternbox}{ButtonModule で学べるパターン}
\begin{itemize}
    \item ModuleTimerを使ったデバウンス処理
    \item エッジ検出（justPressed / justReleased）
    \item 入力モジュールの基本形（ハードウェアから読んでSystemDataに書くだけ）
\end{itemize}
\end{patternbox}

\subsection{Config/Data構造体}

\begin{lstlisting}[style=cppstyle, caption={ButtonModule.h — Config/Data}]
struct ButtonConfig {
    uint8_t       pin;         // ボタンピン
    bool          activeLow;   // true: LOW=押下（プルアップ）
    unsigned long debounceMs;  // デバウンス安定時間 [ms]
};

struct ButtonData {
    bool pressed      = false;  // デバウンス後の押下状態
    bool justPressed  = false;  // 今ループで押下された（立ち上がりエッジ）
    bool justReleased = false;  // 今ループで離された（立ち下がりエッジ）
};
\end{lstlisting}

\begin{pointbox}[Configに動作パラメータを含める]
\texttt{debounceMs}はピンアサインではないが、モジュールの動作特性を決める設定なのでConfigに含める。
ハードウェア環境によって最適なデバウンス時間は異なるため、外部から注入可能にしておく。
\end{pointbox}

\subsection{デバウンスの実装}

\begin{lstlisting}[style=cppstyle, caption={ButtonModule.cpp — updateInput()}]
void ButtonModule::updateInput(SystemData& data) {
    // 生の入力値を読み取り
    bool rawPressed = digitalRead(_config.pin);
    if (_config.activeLow) {
        rawPressed = !rawPressed;
    }

    // 入力値が変化したらデバウンスタイマーをリセット
    if (rawPressed != _lastRaw) {
        _debounceTimer.setTime();
        _lastRaw = rawPressed;
    }

    // デバウンス安定時間を超えたら状態を確定
    if (_debounceTimer.getNowTime() >= _config.debounceMs) {
        _stableState = _lastRaw;
    }

    // エッジ検出
    data.button.pressed      = _stableState;
    data.button.justPressed  = (_stableState && !_prevStable);
    data.button.justReleased = (!_stableState && _prevStable);
    _prevStable = _stableState;
}
\end{lstlisting}

\begin{pointbox}[ModuleTimerによるデバウンス]
仕様書のデバウンスパターンそのまま。
入力が変化するたびに\texttt{setTime()}でタイマーをリセットし、
\texttt{debounceMs}の間安定したら状態を確定する。
\texttt{delay()}は一切使わず、毎ループ非同期で判定している。
\end{pointbox}

\begin{pointbox}[エッジ検出をData構造体に含める]
\texttt{justPressed}と\texttt{justReleased}をData構造体に含めることで、
ロジックフェーズで「ボタンが今押された瞬間」を簡単に検出できる。
\texttt{if (data.button.justPressed)}と書くだけでよく、
ロジック側で前回値を保持する必要がない。
\end{pointbox}

\subsection{ロジックフェーズでの使用例}

\begin{lstlisting}[style=cppstyle, caption={ロジックフェーズでのButtonData参照}]
void applyPattern(SystemData& data) {
    // ボタンが押された瞬間にLEDをトグル
    if (data.button.justPressed) {
        data.led.state = !data.led.state;
    }

    // ボタン押下中にブザーを鳴らす
    if (data.button.pressed) {
        data.buzzer.frequency = 1000;
        data.buzzer.requestPlay = true;
    }
}
\end{lstlisting}

% ─── API仕様 ───────────────────────────────
\subsection{API仕様}


\begin{apibox}{ButtonModule --- ボタン入力モジュール}
\textbf{定義ファイル:} \texttt{lib/ButtonModule/ButtonModule.h / .cpp}\\
\textbf{分類:} 入力モジュール（\texttt{updateInput()}のみオーバーライド）\\
\textbf{対象デバイス:} 汎用タクトスイッチ・プッシュボタン

GPIO入力をデバウンス処理し、押下状態とエッジ（立ち上がり・立ち下がり）を検出するモジュール。
プルアップ接続（activeLow）とプルダウン接続の両方に対応する。
\end{apibox}

\subsubsection{機能概要}

\begin{itemize}
    \item デバウンスフィルタ（\texttt{ModuleTimer}ベース、設定可能な安定時間）
    \item プルアップ（LOW=押下）/ プルダウン（HIGH=押下）の切り替え対応
    \item エッジ検出フラグ（\texttt{justPressed} / \texttt{justReleased}）
    \item 安定状態の保持（\texttt{pressed}）
\end{itemize}

% --- Config ---
\subsubsection{Config構造体}

\begin{funcblock}{ButtonConfig}
ボタンのハードウェア接続とデバウンスパラメータを定義する。

\begin{longtable}{l l l p{5.2cm}}
\toprule
\textbf{メンバ} & \textbf{型} & \textbf{制約・既定値} & \textbf{説明} \\
\midrule
\endhead
\texttt{pin}        & \texttt{uint8\_t}       & ---          & ボタン入力GPIOピン番号 \\
\texttt{activeLow}  & \texttt{bool}           & ---          & \texttt{true}: LOW=押下（プルアップ接続）、\texttt{false}: HIGH=押下（プルダウン接続） \\
\texttt{debounceMs} & \texttt{unsigned long}  & 例: 30       & デバウンス安定時間 [ms]。入力変化後、この時間が経過するまで状態を確定しない \\
\bottomrule
\end{longtable}
\end{funcblock}

% --- Data ---
\subsubsection{Data構造体}

\begin{funcblock}{ButtonData}
デバウンス処理後のボタン状態を保持する。\texttt{justPressed}/\texttt{justReleased}は
1ループ限りのエッジフラグであり、次の\texttt{updateInput()}呼び出しでクリアされる。

\begin{longtable}{l l l p{5.2cm}}
\toprule
\textbf{メンバ} & \textbf{型} & \textbf{初期値} & \textbf{説明} \\
\midrule
\endhead
\texttt{pressed}      & \texttt{bool} & \texttt{false} & デバウンス後の押下状態。押されている間\texttt{true} \\
\texttt{justPressed}  & \texttt{bool} & \texttt{false} & 今ループで押下された瞬間のみ\texttt{true}（立ち上がりエッジ） \\
\texttt{justReleased} & \texttt{bool} & \texttt{false} & 今ループで離された瞬間のみ\texttt{true}（立ち下がりエッジ） \\
\bottomrule
\end{longtable}
\end{funcblock}

% --- メソッド ---
\subsubsection{メソッド}

\begin{funcblock}{ButtonModule(const ButtonConfig\& config)}
\begin{tabularx}{\textwidth}{l X}
\toprule
\textbf{項目} & \textbf{内容} \\
\midrule
引数   & \texttt{config} --- \texttt{ButtonConfig}へのconst参照 \\
動作   & Configを\texttt{\_config}にコピー保持する \\
\bottomrule
\end{tabularx}
\end{funcblock}

\begin{funcblock}{bool init()}
\begin{tabularx}{\textwidth}{l X}
\toprule
\textbf{項目} & \textbf{内容} \\
\midrule
戻り値   & \texttt{bool} --- 常に\texttt{true} \\
動作     & \texttt{activeLow}に応じて\texttt{INPUT\_PULLUP}または\texttt{INPUT}でピンモードを設定。デバウンスタイマーを初期化 \\
\bottomrule
\end{tabularx}
\end{funcblock}

\begin{funcblock}{void updateInput(SystemData\& data)}
\begin{tabularx}{\textwidth}{l X}
\toprule
\textbf{項目} & \textbf{内容} \\
\midrule
書き込み先     & \texttt{data.button} \\
タイミング制御 & 毎ループ実行（タイマーはデバウンス判定に使用） \\
動作           & \texttt{digitalRead()}で生入力を取得し、前回値と比較。変化があればデバウンスタイマーをリセット。タイマー経過後に安定状態を確定し、\texttt{pressed}・\texttt{justPressed}・\texttt{justReleased}を更新 \\
\bottomrule
\end{tabularx}
\end{funcblock}

\begin{notebox}[エッジフラグのライフサイクル]
\texttt{justPressed}と\texttt{justReleased}は\texttt{updateInput()}の先頭で毎回\texttt{false}にリセットされる。
状態変化があったループでのみ\texttt{true}になるため、ロジックフェーズで1回だけ検出できる。
複数箇所で参照しても安全である（同一ループ内では値が変わらない）。
\end{notebox}

% --- 使用例 ---
\subsubsection{使用例}

\begin{lstlisting}[style=cppstyle, caption={ButtonModuleの使用例}]
// ProjectConfig.h
const ButtonConfig BUTTON_CONFIG = {
    .pin        = 0,     // GPIO0: BOOTボタン
    .activeLow  = true,  // LOW=押下（内蔵プルアップ使用）
    .debounceMs = 30,    // デバウンス30ms
};

// ロジックフェーズでの使用
void applyPattern(SystemData& data) {
    // ボタンが押された瞬間にLEDをトグル
    if (data.button.justPressed) {
        data.led.state = !data.led.state;
    }
}
\end{lstlisting}

% BatteryModule
\section{【基本】BatteryModule — ADC + 移動平均フィルタ}

\begin{patternbox}{BatteryModule で学べるパターン}
\begin{itemize}
    \item ModuleTimerによる周期サンプリング
    \item リングバッファを使った移動平均フィルタ
    \item ADC値の電圧変換（分圧回路対応）
    \item isValidフラグによる「データが十分溜まったか」の表現
\end{itemize}
\end{patternbox}

\subsection{Config/Data構造体}

\begin{lstlisting}[style=cppstyle, caption={BatteryModule.h — Config/Data}]
static constexpr uint8_t BATTERY_FILTER_SIZE = 16;

struct BatteryConfig {
    uint8_t       adcPin;
    float         voltageDivider;      // 分圧比（例: 2.0 = 1/2分圧）
    float         adcRefVoltage;       // ADC基準電圧 [V]
    uint16_t      adcResolution;       // ADC分解能（例: ESP32は4095、AVRは1023）
    unsigned long sampleIntervalMs;
    float         lowVoltageThreshold; // 低電圧警告閾値 [V]
};

struct BatteryData {
    float voltage    = 0.0f;   // フィルタ済み電圧 [V]
    float rawVoltage = 0.0f;   // 最新の生電圧 [V]
    bool  isValid    = false;  // サンプルが十分溜まったらtrue
    bool  isLow      = false;  // 低電圧警告フラグ
};
\end{lstlisting}

\begin{pointbox}[Configにハードウェア特性を含める]
\texttt{voltageDivider}（分圧比）や\texttt{adcRefVoltage}（基準電圧）はハードウェア構成に依存する値。
これらをConfigに含めることで、異なる回路構成でも同じモジュールコードを再利用できる。
\end{pointbox}

\subsection{リングバッファと移動平均}

\begin{lstlisting}[style=cppstyle, caption={BatteryModule.cpp — updateInput()}]
void BatteryModule::updateInput(SystemData& data) {
    if (_sampleTimer.getNowTime() < _config.sampleIntervalMs) {
        return;  // 周期未到達なら何もしない（ノンブロッキング）
    }
    _sampleTimer.setTime();

    // ADC読み取り -> 電圧変換
    uint16_t adcRaw = analogRead(_config.adcPin);
    float voltage = _adcToVoltage(adcRaw);

    // リングバッファに追加
    _samples[_sampleIndex] = voltage;
    _sampleIndex = (_sampleIndex + 1) % BATTERY_FILTER_SIZE;
    if (_sampleCount < BATTERY_FILTER_SIZE) {
        _sampleCount++;
    }

    // SystemDataに書き込み
    data.battery.rawVoltage = voltage;
    data.battery.voltage    = _calcAverage();
    data.battery.isValid    = (_sampleCount >= BATTERY_FILTER_SIZE);
    data.battery.isLow      = (data.battery.voltage < _config.lowVoltageThreshold);
}
\end{lstlisting}

\begin{pointbox}[isValidの使い方]
\texttt{isValid}は「移動平均のサンプルが\texttt{BATTERY\_FILTER\_SIZE}に達したか」を示す。
起動直後はサンプルが少なく平均値が不正確なため、ロジックフェーズは\texttt{isValid == true}になるまで
電圧値を信用しないようにできる。
仕様書の「デフォルト値は正常値と区別できる無効値」ルールと同じ考え方。
\end{pointbox}

\begin{pointbox}[周期サンプリングの定石]
\texttt{updateInput()}の冒頭で\texttt{getNowTime() < interval}なら即\texttt{return}するパターンは、
周期実行の定石。ガード節として最初に書くことで、残りのコードのインデントが深くならない。
\end{pointbox}

% ServoModule
\section{【基本】ServoModule — PWM出力パターン}

\begin{patternbox}{ServoModule で学べるパターン}
\begin{itemize}
    \item LEDC（PWM）の初期化と制御
    \item 「変化がある場合のみ更新」による無駄な出力の抑制
    \item currentAngle による出力状態のフィードバック
\end{itemize}
\end{patternbox}

\subsection{Config/Data構造体}

\begin{lstlisting}[style=cppstyle, caption={ServoModule.h — Config/Data}]
struct ServoConfig {
    uint8_t  pin;            // サーボ信号ピン
    uint8_t  pwmChannel;     // LEDCチャネル番号（0-15）
    uint16_t minPulseUs;     // 最小パルス幅 [us]（例: 500）
    uint16_t maxPulseUs;     // 最大パルス幅 [us]（例: 2500）
    uint8_t  defaultAngle;   // 初期角度 [度]（0-180）
};

struct ServoData {
    uint8_t targetAngle  = 90;  // 目標角度（ロジックフェーズが設定）
    uint8_t currentAngle = 90;  // 現在の角度（モジュールが更新）
};
\end{lstlisting}

\begin{pointbox}[targetとcurrentの分離]
\texttt{targetAngle}はロジックフェーズが書き込む「指示値」、
\texttt{currentAngle}はモジュールが更新する「実際の出力値」。
この分離により、ロジックフェーズが「今サーボがどの角度にいるか」を確認できる。
\end{pointbox}

\subsection{updateOutput()の実装}

\begin{lstlisting}[style=cppstyle, caption={ServoModule.cpp — updateOutput()}]
void ServoModule::updateOutput(SystemData& data) {
    uint8_t target = constrain(data.servo.targetAngle, 0, 180);

    // 変化がある場合のみPWM更新
    if (target != data.servo.currentAngle) {
        ledcWrite(_config.pwmChannel, _angleToDuty(target));
        data.servo.currentAngle = target;
    }
}
\end{lstlisting}

\begin{pointbox}[変化検出パターン]
毎ループPWMを書き込むのではなく、\texttt{target != current}の場合のみ更新する。
同じ値を毎回書き込んでも動作に問題はないが、
不要なハードウェアアクセスを避けるのは良い習慣。
特にI2C/SPI経由のデバイスでは通信コストの削減になる。
\end{pointbox}

\begin{pointbox}[constrainによる安全性]
\texttt{constrain()}でサーボの物理的な可動範囲に制限している。
ロジックフェーズが誤って範囲外の値を設定しても、モジュール側で安全に丸める。
入力バリデーションはモジュールの責務。
\end{pointbox}

% ─── API仕様 ───────────────────────────────
\subsection{API仕様}


\begin{apibox}{ServoModule --- サーボモーター制御モジュール}
\textbf{定義ファイル:} \texttt{lib/ServoModule/ServoModule.h / .cpp}\\
\textbf{分類:} 出力モジュール（\texttt{updateOutput()}のみオーバーライド）\\
\textbf{対象デバイス:} PWMサーボモーター（SG90等）

LEDCによるPWMパルス幅制御でサーボモーターの角度を設定するモジュール。
角度（0--180$^{\circ}$）からパルス幅（$\mu$s）への変換を内部で行う。
\end{apibox}

\subsubsection{機能概要}

\begin{itemize}
    \item 角度指定（0--180$^{\circ}$）によるサーボ制御
    \item パルス幅範囲のカスタマイズ（\texttt{minPulseUs} / \texttt{maxPulseUs}）
    \item 角度変化時のみPWM更新（不要な書き込みを抑制）
    \item 初期角度の設定（\texttt{defaultAngle}）
\end{itemize}

% --- Config ---
\subsubsection{Config構造体}

\begin{funcblock}{ServoConfig}
サーボのPWMパラメータを定義する。

\begin{longtable}{l l l p{4.5cm}}
\toprule
\textbf{メンバ} & \textbf{型} & \textbf{制約・既定値} & \textbf{説明} \\
\midrule
\endhead
\texttt{pin}          & \texttt{uint8\_t}  & ---       & サーボ信号出力GPIOピン番号 \\
\texttt{pwmChannel}   & \texttt{uint8\_t}  & 0--15     & LEDCチャネル番号。他モジュールと重複しないこと \\
\texttt{minPulseUs}   & \texttt{uint16\_t} & 例: 500   & 0$^{\circ}$ 時のパルス幅 [$\mu$s] \\
\texttt{maxPulseUs}   & \texttt{uint16\_t} & 例: 2500  & 180$^{\circ}$ 時のパルス幅 [$\mu$s] \\
\texttt{defaultAngle} & \texttt{uint8\_t}  & 例: 90    & 初期角度 [度]（0--180） \\
\bottomrule
\end{longtable}
\end{funcblock}

% --- Data ---
\subsubsection{Data構造体}

\begin{funcblock}{ServoData}
サーボの目標角度と現在角度を保持する。

\begin{longtable}{l l l p{5.0cm}}
\toprule
\textbf{メンバ} & \textbf{型} & \textbf{初期値} & \textbf{説明} \\
\midrule
\endhead
\texttt{targetAngle}  & \texttt{uint8\_t} & \texttt{90} & 目標角度 [度]（0--180）。ロジックフェーズで設定 \\
\texttt{currentAngle} & \texttt{uint8\_t} & \texttt{90} & 現在のPWM出力に対応する角度 [度]（0--180） \\
\bottomrule
\end{longtable}
\end{funcblock}

% --- メソッド ---
\subsubsection{メソッド}

\begin{funcblock}{ServoModule(const ServoConfig\& config)}
\begin{tabularx}{\textwidth}{l X}
\toprule
\textbf{項目} & \textbf{内容} \\
\midrule
引数   & \texttt{config} --- \texttt{ServoConfig}へのconst参照 \\
動作   & Configを\texttt{\_config}にコピー保持する \\
\bottomrule
\end{tabularx}
\end{funcblock}

\begin{funcblock}{bool init()}
\begin{tabularx}{\textwidth}{l X}
\toprule
\textbf{項目} & \textbf{内容} \\
\midrule
戻り値   & \texttt{bool} --- 常に\texttt{true} \\
動作     & LEDCチャネルを50Hz（サーボ標準）で初期設定し、\texttt{defaultAngle}に対応するパルス幅を出力する \\
\bottomrule
\end{tabularx}
\end{funcblock}

\begin{funcblock}{void updateOutput(SystemData\& data)}
\begin{tabularx}{\textwidth}{l X}
\toprule
\textbf{項目} & \textbf{内容} \\
\midrule
読み取り元     & \texttt{data.servo} \\
タイミング制御 & 毎ループ実行（ただし角度変化なしの場合はPWM更新をスキップ） \\
動作           & \texttt{targetAngle}と\texttt{currentAngle}を比較し、変化があれば角度をパルス幅に変換してPWMデューティを更新。\texttt{currentAngle}を更新 \\
\bottomrule
\end{tabularx}
\end{funcblock}

\begin{notebox}[角度→パルス幅変換]
角度からパルス幅への変換は以下の線形補間で行う:

\vspace{0.3em}
$$pulseWidth = minPulseUs + \frac{angle}{180} \times (maxPulseUs - minPulseUs)$$
\vspace{0.3em}

\noindent SG90の場合（500--2500$\mu$s）: 90$^{\circ}$ $\rightarrow$ 1500$\mu$s（中央）
\end{notebox}

% --- 使用例 ---
\subsubsection{使用例}

\begin{lstlisting}[style=cppstyle, caption={ServoModuleの使用例}]
// ProjectConfig.h
const ServoConfig SERVO_CONFIG = {
    .pin          = 1,     // GPIO1: サーボ信号
    .pwmChannel   = 0,     // LEDCチャネル0
    .minPulseUs   = 500,   // SG90: 500us
    .maxPulseUs   = 2500,  // SG90: 2500us
    .defaultAngle = 90,    // 初期角度 90度（中央）
};

// ロジックフェーズでの使用
void applyPattern(SystemData& data) {
    // IMUの傾きに応じてサーボ角度を設定
    if (data.mpu.isValid) {
        float tilt = atan2(data.mpu.accelX, data.mpu.accelZ);
        data.servo.targetAngle = constrain(
            map(tilt * 57.3, -45, 45, 0, 180), 0, 180
        );
    }
}
\end{lstlisting}

% BuzzerModule
\section{【基本】BuzzerModule — deinit()パターン}

\begin{patternbox}{BuzzerModule で学べるパターン}
\begin{itemize}
    \item deinit()によるPWMリソースの明示的解放
    \item リクエスト方式による出力制御（requestPlayフラグ）
    \item ModuleTimerによる時間指定再生の自動停止
\end{itemize}
\end{patternbox}

\subsection{Data構造体のリクエストパターン}

\begin{lstlisting}[style=cppstyle, caption={BuzzerModule.h — Data}]
struct BuzzerData {
    uint16_t frequency = 0;       // 再生周波数 [Hz]（0=消音）
    unsigned long durationMs = 0; // 再生時間 [ms]（0=無制限）
    bool     requestPlay = false; // ロジックフェーズからの再生リクエスト
};
\end{lstlisting}

\begin{pointbox}[リクエストフラグパターン]
\texttt{requestPlay}は「ロジックフェーズからの指示を1回だけ処理する」ためのフラグ。
\texttt{updateOutput()}が処理後に\texttt{false}に戻すため、次ループでは再度処理しない。
LedModuleのように毎ループ同じ状態を出力するのではなく、
「イベント的に1回だけ何かをする」場合にこのパターンを使う。
\end{pointbox}

\subsection{updateOutput()の実装}

\begin{lstlisting}[style=cppstyle, caption={BuzzerModule.cpp — updateOutput()}]
void BuzzerModule::updateOutput(SystemData& data) {
    // 再生リクエスト処理
    if (data.buzzer.requestPlay) {
        data.buzzer.requestPlay = false;  // 処理済みフラグクリア

        if (data.buzzer.frequency > 0) {
            _toneOn(data.buzzer.frequency);
            if (data.buzzer.durationMs > 0) {
                _durationTimer.setTime();
                _timed = true;
            } else {
                _timed = false;  // 無制限再生
            }
        } else {
            _toneOff();
        }
    }

    // 時間指定ありの場合、経過したら自動消音
    if (_playing && _timed) {
        if (_durationTimer.getNowTime() >= data.buzzer.durationMs) {
            _toneOff();
        }
    }
}
\end{lstlisting}

\subsection{deinit()の実装}

\begin{lstlisting}[style=cppstyle, caption={BuzzerModule.cpp — deinit()}]
void BuzzerModule::deinit() {
    _toneOff();
    ledcDetachPin(_config.pin);  // PWMチャネルとピンの紐付けを解除
    Serial.println("[Buzzer] deinit OK");
}
\end{lstlisting}

\begin{pointbox}[deinit()の実装が必要なケース]
\texttt{IModule::deinit()}はデフォルトで空実装。
BuzzerModuleのように\textbf{ハードウェアリソース（PWMチャネル）を占有する}モジュールでは、
スリープ前やモジュール停止時にリソースを明示的に解放する必要がある。\\
\texttt{ledcDetachPin()}を呼ばないと、PWMチャネルが占有されたままスリープに入り、
復帰後に同じチャネルを使えなくなる可能性がある。
\end{pointbox}

\begin{notebox}[deinit()が不要なモジュール]
LedModuleやButtonModuleは\texttt{pinMode()}で設定したGPIOを使うだけで、
特別なリソースを占有しない。こうしたモジュールでは\texttt{deinit()}をオーバーライドしなくてよい。
\end{notebox}

% Mpu6500Module
\section{【中級】Mpu6500Module — I2Cバスパターン}

\begin{patternbox}{Mpu6500Module で学べるパターン}
\begin{itemize}
    \item 通信バスパターン（コンストラクタでバスポインタを受け取る）
    \item ヘッダーでのバスクラス前方宣言（\texttt{<Wire.h>}依存を.cppに閉じ込める）
    \item WHO\_AM\_I による init() でのデバイス識別
    \item 周期サンプリング + バースト読み取り
\end{itemize}
\end{patternbox}

\subsection{バスポインタの受け取り}

\begin{lstlisting}[style=cppstyle, caption={Mpu6500Module.h — バスの前方宣言とコンストラクタ}]
// TwoWireはポインタのみ使用するため前方宣言で十分
class TwoWire;

class Mpu6500Module : public IModule {
private:
    Mpu6500Config _config;
    TwoWire*      _wire;       // コンストラクタで受け取るバスポインタ
    ModuleTimer   _sampleTimer;
public:
    explicit Mpu6500Module(const Mpu6500Config& config, TwoWire* wire);
    // ...
};
\end{lstlisting}

\begin{pointbox}[ヘッダーでの前方宣言]
\texttt{class TwoWire;}と前方宣言することで、ヘッダーに\texttt{<Wire.h>}をインクルードする必要がない。
ポインタや参照の宣言には前方宣言だけで十分。
これにより、\texttt{Mpu6500Module.h}をインクルードするファイルに\texttt{Wire.h}の依存が波及しない。
\end{pointbox}

\begin{pointbox}[Configにバスポインタを含めない理由]
\texttt{Mpu6500Config}にはI2Cアドレスとピン情報だけを含め、\texttt{TwoWire*}は含めない。
バスポインタをConfigに含めると、\texttt{ProjectConfig.h}でバスインスタンスの\texttt{extern}宣言が
必要になり、インクルード順序の制約が生じる。
コンストラクタの別引数にすることで、Configは純粋な設定値のみを持つ。
\end{pointbox}

\subsection{main.cppでのバス管理}

\begin{lstlisting}[style=cppstyle, caption={main.cpp — バスの生成と初期化}]
// バスインスタンス（グローバルスコープで生成）
static TwoWire mpuWire = TwoWire(0);

// モジュール生成（Config + バスポインタ）
Mpu6500Module mpu6500Module(MPU6500_CONFIG, &mpuWire);

void setup() {
    // バス初期化は全モジュールのinit()より前に実行
    mpuWire.begin(MPU6500_CONFIG.sdaPin, MPU6500_CONFIG.sclPin);
    // ...
    mpu6500Module.init();  // init()ではbus.begin()を呼ばない
}
\end{lstlisting}

\begin{pointbox}[init()でbus.begin()を呼ばない]
バスの初期化は\texttt{main.cpp}の\texttt{setup()}で一元管理する。
モジュールの\texttt{init()}ではデバイス固有の初期化（WHO\_AM\_Iの確認、レジスタ設定等）のみ行う。
同じバスを複数モジュールが共有する場合、各モジュールが\texttt{begin()}すると二重初期化になる。
\end{pointbox}

% TouchModule
\section{【中級】TouchModule — バスオブジェクト共有パターン}

\begin{patternbox}{TouchModule で学べるパターン}
\begin{itemize}
    \item 外部ライブラリのドライバオブジェクト共有（TFT\_eSPIをTftModuleと共有）
    \item 入出力分離の原則（タッチ入力とLCD出力を別モジュールに分ける）
    \item 空のConfig構造体（設定がplatformio.iniで管理される場合）
\end{itemize}
\end{patternbox}

\subsection{ドライバオブジェクトの共有}

\begin{lstlisting}[style=cppstyle, caption={TouchModule.h — TFT\_eSPIの前方宣言}]
class TFT_eSPI;  // TftModuleと共有するドライバ

class TouchModule : public IModule {
private:
    TouchConfig _config;
    TFT_eSPI*   _tft;  // コンストラクタで受け取る共有ドライバ
public:
    TouchModule(const TouchConfig& config, TFT_eSPI* tft);
    // ...
};
\end{lstlisting}

\begin{lstlisting}[style=cppstyle, caption={main.cpp — ドライバの共有}]
static TFT_eSPI tftDriver;  // 1つのインスタンスを生成

// 入力と出力の両モジュールが同じドライバを共有
TouchModule touchModule(TOUCH_CONFIG, &tftDriver);
TftModule   tftModule(TFT_CONFIG, &tftDriver);
\end{lstlisting}

\begin{pointbox}[通信バスパターンの応用]
I2C/SPIのバスインスタンスだけでなく、外部ライブラリのドライバオブジェクトも同じパターンで共有できる。
TFT\_eSPIはSPIバスの抽象化であり、タッチ機能もSPIを共有するため、
1つのドライバインスタンスをTouchModuleとTftModuleの両方に渡す。
\end{pointbox}

\begin{pointbox}[入出力分離の実例]
TFTディスプレイはタッチ入力とLCD出力の両方を持つが、
\textbf{TouchModule（入力）とTftModule（出力）に分離}している。
入力フェーズでタッチ座標を読み取り、出力フェーズで画面を描画する。
これにより各モジュールの責務が「入力のみ」「出力のみ」に限定される。
\end{pointbox}

% TftModule
\section{【中級】TftModule — SPI共有 + 複合デバイスの分離}

\begin{patternbox}{TftModule で学べるパターン}
\begin{itemize}
    \item SPI共有ドライバを使った画面描画
    \item ModuleTimerによるFPS制御（更新頻度の制限）
    \item テキスト表示をData構造体経由で受け渡す方法
\end{itemize}
\end{patternbox}

\subsection{Data構造体のテキストバッファ}

\begin{lstlisting}[style=cppstyle, caption={TftModule.h — Data}]
struct TftData {
    bool touchPressed = false;
    int16_t touchX = 0, touchY = 0;
    char line1[48] = "";  // ロジックフェーズが書き込む表示テキスト
    char line2[48] = "";
    char line3[48] = "";
    char line4[48] = "";
    char line5[48] = "";
};
\end{lstlisting}

\begin{pointbox}[出力モジュールは「何を表示するか」を決めない]
\texttt{line1}〜\texttt{line5}はロジックフェーズが\texttt{snprintf()}で書き込む文字列。
TftModuleの\texttt{update()}はそれを画面に描画するだけで、
表示内容の判断は一切行わない。出力モジュールの責務を守っている。
\end{pointbox}

\subsection{ロジックフェーズでの表示内容設定}

\begin{lstlisting}[style=cppstyle, caption={main.cpp — applyPattern()でのTFT表示設定}]
void applyPattern(SystemData& data) {
    // MPU-6500データをTFT表示用にフォーマット
    if (data.mpu.isValid) {
        snprintf(data.tft.line1, sizeof(data.tft.line1),
                 "Ax:%.2f Ay:%.2f Az:%.2f",
                 data.mpu.accelX, data.mpu.accelY, data.mpu.accelZ);
    }
    // カメラ情報
    if (data.camera.isValid && data.camera.frameReady) {
        snprintf(data.tft.line4, sizeof(data.tft.line4),
                 "CAM: %dx%d (%d B)",
                 data.camera.width, data.camera.height,
                 (int)data.camera.frameSize);
    }
}
\end{lstlisting}

\begin{pointbox}[入力→ロジック→出力のデータフロー]
入力フェーズでMpu6500Moduleが\texttt{data.mpu}に値を書き込み、
ロジックフェーズでそれを文字列に変換して\texttt{data.tft.line1}に書き込み、
出力フェーズでTftModuleが画面に描画する。
3フェーズモデルの典型的なデータフローの実例。
\end{pointbox}

% CameraModule
\section{【中級】CameraModule — PSRAMバッファ + フレーム管理}

\begin{patternbox}{CameraModule で学べるパターン}
\begin{itemize}
    \item PSRAMへの大容量バッファ配置
    \item フレームバッファのライフサイクル管理（取得→使用→解放）
    \item Data構造体にメタデータのみ格納し、生データはモジュール内に保持
    \item 多数のピンをConfig構造体で一元管理
\end{itemize}
\end{patternbox}

\subsection{Configの大量ピン管理}

\begin{lstlisting}[style=cppstyle, caption={CameraModule.h — Config（多数のピン）}]
struct CameraConfig {
    int8_t pwdnPin, resetPin;           // 電源/リセット
    int8_t xclkPin, siodPin, siocPin;   // クロック/SCCB
    int8_t d0Pin, d1Pin, d2Pin, d3Pin;  // DVPデータバス
    int8_t d4Pin, d5Pin, d6Pin, d7Pin;
    int8_t vsyncPin, hrefPin, pclkPin;  // 同期
    int     xclkFreqHz;
    uint8_t frameSize, pixFormat;
    int     jpegQuality;
    size_t  fbCount;
};
\end{lstlisting}

\begin{pointbox}[Config構造体の威力]
カメラモジュールはピンが16本以上ある。
コンストラクタの引数で16個のピンを渡すと順序を間違える危険があるが、
Config構造体なら\texttt{.d0Pin = 11}のように名前で指定できる。
Config構造体のピンアサインが\texttt{ProjectConfig.h}に集約されるため、
回路設計者はこのファイルだけでピンの確認・変更ができる。
\end{pointbox}

\subsection{フレームバッファの管理}

\begin{lstlisting}[style=cppstyle, caption={CameraModule — フレームバッファ管理}]
// updateInput()はメタデータだけをSystemDataに書き込む
void CameraModule::updateInput(SystemData& data) {
    // 保持中のフレームがあれば先に解放
    releaseFrame();

    camera_fb_t* fb = esp_camera_fb_get();
    if (fb) {
        data.camera.frameReady = true;
        data.camera.width = fb->width;
        data.camera.height = fb->height;
        data.camera.frameSize = fb->len;
        _frameBuffer = fb;  // void*で保持（ヘッダ依存回避）
        _frameHeld = true;
    }
}

// ロジックフェーズから生データにアクセスする場合
const uint8_t* CameraModule::getFrameBuffer() const {
    if (!_frameHeld || !_frameBuffer) return nullptr;
    return ((camera_fb_t*)_frameBuffer)->buf;
}
void CameraModule::releaseFrame() {
    if (_frameHeld && _frameBuffer) {
        esp_camera_fb_return((camera_fb_t*)_frameBuffer);
        _frameBuffer = nullptr;
        _frameHeld = false;
    }
}
\end{lstlisting}

\begin{pointbox}[Data構造体に生データを入れない]
\texttt{camera\_fb\_t*}（フレームバッファポインタ）はSystemDataに格納しない。
SystemDataに入れると、他モジュールがフレームの解放タイミングを制御できてしまう。
代わりに\texttt{getFrameBuffer()} / \texttt{releaseFrame()}でアクセスを制御し、
\texttt{CameraData}にはメタデータ（サイズ、解像度等）のみを格納する。
\end{pointbox}

\begin{notebox}[PSRAMの自動配置]
\texttt{fbCount >= 2}の場合、esp\_cameraライブラリが自動的にPSRAMにバッファを配置する。
Config構造体の\texttt{fbCount}で制御でき、モジュール内部でPSRAM判定のコードを書く必要はない。
\end{notebox}

% ─── API仕様 ───────────────────────────────
\subsection{API仕様}


\begin{apibox}{CameraModule --- OV5640カメラモジュール}
\textbf{定義ファイル:} \texttt{lib/CameraModule/CameraModule.h / .cpp}\\
\textbf{分類:} 入力モジュール（\texttt{updateInput()}のみオーバーライド）\\
\textbf{対象デバイス:} OV5640（ESP32-S3 DVP接続）

ESP32-S3のDVPインターフェース経由でOV5640カメラからフレームを取得するモジュール。
\texttt{esp\_camera}ライブラリを使用し、フレームバッファの取得・解放を管理する。
フレームの生データは\texttt{SystemData}には格納せず、専用メソッドでアクセスする。
\end{apibox}

\subsubsection{機能概要}

\begin{itemize}
    \item DVP 16ピンインターフェースによるカメラ接続
    \item JPEG/RAWフォーマット対応（\texttt{pixFormat}で選択）
    \item フレームバッファの取得・解放API（\texttt{getFrameBuffer()} / \texttt{releaseFrame()}）
    \item フレームメタデータ（解像度・サイズ）の\texttt{SystemData}への書き込み
    \item \texttt{deinit()}によるカメラリソース解放
\end{itemize}

% --- Config ---
\subsubsection{Config構造体}

\begin{funcblock}{CameraConfig}
DVPピンアサインとカメラ動作パラメータを定義する。

\begin{longtable}{l l l p{4.2cm}}
\toprule
\textbf{メンバ} & \textbf{型} & \textbf{制約・既定値} & \textbf{説明} \\
\midrule
\endfirsthead
\toprule
\textbf{メンバ} & \textbf{型} & \textbf{制約・既定値} & \textbf{説明} \\
\midrule
\endhead
\midrule
\multicolumn{4}{r}{\small 次ページに続く} \\
\endfoot
\bottomrule
\endlastfoot
\multicolumn{4}{l}{\textbf{--- DVPピンアサイン ---}} \\
\texttt{pwdnPin}    & \texttt{int8\_t} & -1で未使用         & POWER DOWNピン \\
\texttt{resetPin}   & \texttt{int8\_t} & -1で未使用         & ハードウェアRESETピン \\
\texttt{xclkPin}    & \texttt{int8\_t} & ---                & XCLKクロック出力ピン \\
\texttt{siodPin}    & \texttt{int8\_t} & ---                & SCCB SDA（カメラ設定用I2C） \\
\texttt{siocPin}    & \texttt{int8\_t} & ---                & SCCB SCL \\
\texttt{d0Pin}--\texttt{d7Pin} & \texttt{int8\_t} & --- & データバス D0--D7（8ピン） \\
\texttt{vsyncPin}   & \texttt{int8\_t} & ---                & 垂直同期ピン \\
\texttt{hrefPin}    & \texttt{int8\_t} & ---                & 水平参照ピン \\
\texttt{pclkPin}    & \texttt{int8\_t} & ---                & ピクセルクロックピン \\
\midrule
\multicolumn{4}{l}{\textbf{--- カメラ動作設定 ---}} \\
\texttt{xclkFreqHz}  & \texttt{int}     & 通常: 20000000   & XCLKクロック周波数 [Hz] \\
\texttt{frameSize}   & \texttt{uint8\_t} & 例: 5 (QVGA)    & フレームサイズ（\texttt{framesize\_t}の値） \\
\texttt{pixFormat}   & \texttt{uint8\_t} & 例: 4 (JPEG)    & ピクセルフォーマット（\texttt{pixformat\_t}の値） \\
\texttt{jpegQuality} & \texttt{int}     & 0--63（低い=高画質）& JPEG圧縮品質 \\
\texttt{fbCount}     & \texttt{size\_t} & 推奨: 2          & フレームバッファ数 \\
\end{longtable}
\end{funcblock}

% --- Data ---
\subsubsection{Data構造体}

\begin{funcblock}{CameraData}
フレームのメタデータを保持する。フレームの生データ（ピクセルバッファ）は含まない。

\begin{longtable}{l l l p{4.8cm}}
\toprule
\textbf{メンバ} & \textbf{型} & \textbf{初期値} & \textbf{説明} \\
\midrule
\endhead
\texttt{isValid}    & \texttt{bool}     & \texttt{false} & カメラ初期化が成功していれば\texttt{true} \\
\texttt{frameReady} & \texttt{bool}     & \texttt{false} & 新しいフレームが取得済みであれば\texttt{true} \\
\texttt{width}      & \texttt{uint16\_t} & \texttt{0}    & フレーム幅 [px] \\
\texttt{height}     & \texttt{uint16\_t} & \texttt{0}    & フレーム高さ [px] \\
\texttt{frameSize}  & \texttt{size\_t}  & \texttt{0}     & フレームデータサイズ [bytes] \\
\bottomrule
\end{longtable}
\end{funcblock}

\begin{cautionbox}[フレームデータのアクセス方法]
フレームの生データ（ピクセルバッファ）は大容量のためSystemDataには格納しない。
ロジックフェーズで\texttt{getFrameBuffer()}を呼び出してポインタを取得し、
使用後は必ず\texttt{releaseFrame()}で解放すること。
解放しないとフレームバッファが枯渇し、次のフレーム取得がブロックされる。
\end{cautionbox}

% --- メソッド ---
\subsubsection{メソッド}

\begin{funcblock}{CameraModule(const CameraConfig\& config)}
\begin{tabularx}{\textwidth}{l X}
\toprule
\textbf{項目} & \textbf{内容} \\
\midrule
引数   & \texttt{config} --- \texttt{CameraConfig}へのconst参照 \\
動作   & Configを\texttt{\_config}にコピー保持する \\
\bottomrule
\end{tabularx}
\end{funcblock}

\begin{funcblock}{bool init()}
\begin{tabularx}{\textwidth}{l X}
\toprule
\textbf{項目} & \textbf{内容} \\
\midrule
戻り値     & \texttt{bool} --- カメラ初期化成功で\texttt{true}、失敗で\texttt{false} \\
動作       & Configからesp\_cameraの設定構造体を構築し、\texttt{esp\_camera\_init()}を呼び出す \\
失敗条件   & カメラ未接続、ピン設定誤り、PSRAMメモリ不足等 \\
\bottomrule
\end{tabularx}
\end{funcblock}

\begin{funcblock}{void updateInput(SystemData\& data)}
\begin{tabularx}{\textwidth}{l X}
\toprule
\textbf{項目} & \textbf{内容} \\
\midrule
書き込み先     & \texttt{data.camera} \\
タイミング制御 & 毎ループ実行 \\
動作           & 前回のフレームが保持中であれば解放し、新しいフレームを\texttt{esp\_camera\_fb\_get()}で取得。メタデータ（幅、高さ、サイズ）を\texttt{data.camera}に書き込む \\
\bottomrule
\end{tabularx}
\end{funcblock}

\begin{funcblock}{const uint8\_t* getFrameBuffer() const}
\begin{tabularx}{\textwidth}{l X}
\toprule
\textbf{項目} & \textbf{内容} \\
\midrule
戻り値   & フレームデータへのポインタ。フレーム未取得時は\texttt{nullptr} \\
用途     & ロジックフェーズでフレームデータにアクセスする（BLE送信、SD保存、画像処理等） \\
有効期間 & 次の\texttt{updateInput()}呼び出しまで、または\texttt{releaseFrame()}呼び出しまで \\
\bottomrule
\end{tabularx}
\end{funcblock}

\begin{funcblock}{void releaseFrame()}
\begin{tabularx}{\textwidth}{l X}
\toprule
\textbf{項目} & \textbf{内容} \\
\midrule
動作   & 保持中のフレームバッファを\texttt{esp\_camera\_fb\_return()}で解放する。フレーム未保持時は何もしない \\
\bottomrule
\end{tabularx}
\end{funcblock}

\begin{funcblock}{void deinit()}
\begin{tabularx}{\textwidth}{l X}
\toprule
\textbf{項目} & \textbf{内容} \\
\midrule
動作   & 保持中のフレームを解放し、\texttt{esp\_camera\_deinit()}でカメラリソースを解放する \\
\bottomrule
\end{tabularx}
\end{funcblock}

\begin{cautionbox}[PSRAMの必要性]
\texttt{fbCount}を2以上に設定する場合、PSRAMが必要である。
PSRAMがない環境では\texttt{fbCount = 1}に制限し、フレームサイズも小さくする必要がある。
\texttt{platformio.ini}で\texttt{board\_build.psram = enabled}を設定すること。
\end{cautionbox}

% --- 使用例 ---
\subsubsection{使用例}

\begin{lstlisting}[style=cppstyle, caption={CameraModuleの使用例}]
// ProjectConfig.h
const CameraConfig CAMERA_CONFIG = {
    .pwdnPin = -1, .resetPin = -1, .xclkPin = 15,
    .siodPin = 4,  .siocPin = 5,
    .d0Pin = 11, .d1Pin = 9, .d2Pin = 8, .d3Pin = 10,
    .d4Pin = 12, .d5Pin = 18, .d6Pin = 17, .d7Pin = 16,
    .vsyncPin = 6, .hrefPin = 7, .pclkPin = 13,
    .xclkFreqHz  = 20000000,
    .frameSize   = 5,   // QVGA (320x240)
    .pixFormat   = 4,   // JPEG
    .jpegQuality = 12,
    .fbCount     = 2,   // ダブルバッファ（PSRAM必須）
};

// ロジックフェーズでの使用
void applyPattern(SystemData& data) {
    if (data.camera.frameReady) {
        const uint8_t* buf = camera.getFrameBuffer();
        size_t len = data.camera.frameSize;
        // BLE送信やSD保存などの処理
        camera.releaseFrame();
    }
}
\end{lstlisting}

% DriveMotorModule
\section{【応用】DriveMotorModule — PWM 2ピン制御 + 統合部品}

\begin{patternbox}{DriveMotorModule で学べるパターン}
\begin{itemize}
    \item Hブリッジによる方向制御（IN1/IN2 + PWM）
    \item 逆転保護（ブレーキ自動挿入）
    \item drive()メソッドによるSystemData非経由の直接制御
    \item 統合モジュールの「内部部品」としての使い方
\end{itemize}
\end{patternbox}

\subsection{2つの使い方}

\texttt{DriveMotorModule}は2つのモードで使える。

\begin{enumerate}
    \item \textbf{単体使用} — \texttt{outputModules[]}に登録し、\texttt{SystemData}経由で制御
    \item \textbf{統合モジュールの内部部品} — \texttt{ChassisModule}等が内部で\texttt{drive()}を呼ぶ
\end{enumerate}

\begin{lstlisting}[style=cppstyle, caption={DriveMotorModule.h — drive()メソッド}]
class DriveMotorModule : public IModule {
public:
    // SystemData経由の制御（単体使用時）
    void update(SystemData& data) override;

    // SystemData非経由の直接制御（統合モジュールから使用）
    void drive(float power);
    void brake();
};
\end{lstlisting}

\begin{pointbox}[update()とdrive()の使い分け]
\texttt{update()}はSystemDataの\texttt{DriveMotorData.power}を読んで\texttt{drive()}を呼ぶラッパー。
統合モジュールの内部部品として使う場合は\texttt{drive()}を直接呼び、
\texttt{update()}は空実装（呼ばれない）となる。
この設計により、同じモジュールコードを単体でも統合でも使い回せる。
\end{pointbox}

\subsection{逆転保護}

\begin{lstlisting}[style=cppstyle, caption={DriveMotorModule.cpp — 逆転保護}]
void DriveMotorModule::drive(float power) {
    Direction newDir = (power > 0) ? FORWARD
                     : (power < 0) ? REVERSE : STOP;

    // 正転<->逆転の直接切り替え時、1フレーム分のブレーキを挿入
    if ((_lastDir == FORWARD && newDir == REVERSE) ||
        (_lastDir == REVERSE && newDir == FORWARD)) {
        brake();
        _lastDir = STOP;
        return;  // 次のupdate()で実際の方向転換が行われる
    }
    // ... 実際のPWM制御 ...
}
\end{lstlisting}

\begin{pointbox}[1フレーム遅延によるモーター保護]
正転→逆転の直接切り替えはモーターやドライバICに大きな負荷がかかる。
1フレーム分（1ループ分）のブレーキを挿入することで、物理的に安全な方向転換を実現する。
3フェーズモデルの「毎ループ1回だけ呼ばれる」性質を利用した保護パターン。
\end{pointbox}

% ChassisModule
\section{【応用】ChassisModule — 機能統合パターン}

\begin{patternbox}{ChassisModule で学べるパターン}
\begin{itemize}
    \item 複数モジュールを内部に保持する「機能統合系モジュール」
    \item 内部モジュールのinit()/updateOutput()/deinit()を親が管理
    \item ベクトル合成による高レベル制御（3軸入力→4輪出力）
    \item 内部モジュールのSystemData登録が不要なケース
\end{itemize}
\end{patternbox}

\subsection{内部モジュールの保持}

\begin{lstlisting}[style=cppstyle, caption={ChassisModule.h — 内部にDriveMotorModuleを4つ保持}]
struct ChassisConfig {
    DriveMotorConfig motors[4];     // 4輪分のモーターConfig
    WheelPattern     wheelPattern;
    float            maxOutputPercent;
};

class ChassisModule : public IModule {
private:
    ChassisConfig      _config;
    DriveMotorModule   _motors[4];  // 内部にモジュールを直接保持
public:
    ChassisModule(const ChassisConfig& config);
    bool init() override;
    void updateOutput(SystemData& data) override;
    void deinit() override;
};
\end{lstlisting}

\begin{pointbox}[内部モジュールはIModule配列に登録しない]
\texttt{\_motors[4]}は\texttt{ChassisModule}のメンバ変数であり、
\texttt{main.cpp}の\texttt{outputModules[]}には登録しない。
\texttt{SystemData}にも\texttt{DriveMotorData}を含める必要がない。
外部からは\texttt{ChassisModule}1つだけが見え、3軸入力だけで4輪を制御できる。
\end{pointbox}

\subsection{親モジュールからの制御委譲}

\begin{lstlisting}[style=cppstyle, caption={ChassisModule.cpp — init()/deinit()で内部モジュールを管理}]
ChassisModule::ChassisModule(const ChassisConfig& config)
    : _config(config),
      _motors{
          DriveMotorModule(config.motors[0]),
          DriveMotorModule(config.motors[1]),
          DriveMotorModule(config.motors[2]),
          DriveMotorModule(config.motors[3]),
      } {}
\end{lstlisting}

\begin{notebox}[コンストラクタの初期化リスト — 初学者向け解説]
コロン（\texttt{:}）の後に書かれている部分が「初期化リスト」です。ここでは以下の処理を行っています。

\begin{enumerate}
    \item \texttt{\_config(config)} — Config構造体をコピー保存
    \item \texttt{\_motors\{...\}} — DriveMotorModuleの配列を、各モーターのConfigで初期化
\end{enumerate}

配列の初期化は波括弧 \texttt{\{\}} で囲みます。中の \texttt{DriveMotorModule(config.motors[0])} は「config.motors[0]の設定でDriveMotorModuleを作る」という意味です。

\vspace{2mm}
これは「コンストラクタの中で4つの子モジュールを一気に作る」処理です。
初期化リスト構文に慣れないうちは「コンストラクタ引数を受け取って、メンバ変数に渡している」とだけ理解すれば十分です。
\end{notebox}

\begin{lstlisting}[style=cppstyle, caption={ChassisModule.cpp — init()/deinit()（続き）}]

bool ChassisModule::init() {
    for (int i = 0; i < 4; i++) {
        if (!_motors[i].init()) return false;
    }
    return true;
}

void ChassisModule::deinit() {
    for (int i = 0; i < 4; i++) {
        _motors[i].deinit();
    }
}
\end{lstlisting}

\begin{lstlisting}[style=cppstyle, caption={ChassisModule.cpp — updateOutput()でベクトル合成→各モーターをdrive()}]
void ChassisModule::updateOutput(SystemData& data) {
    float outputs[4];
    // ベクトル合成（ホイールパターンに応じた方向行列）
    calcWheelOutputs(data.chassis, outputs);
    // リミッタ + スケーリング
    applyLimiter(outputs, _config.maxOutputPercent);
    // 各内部モジュールのdrive()を直接呼ぶ
    for (int i = 0; i < 4; i++) {
        _motors[i].drive(outputs[i]);
        data.chassis.motorOutputs[i] = outputs[i];
    }
}
\end{lstlisting}

\begin{pointbox}[機能統合の原則]
仕様書の「モジュール間の依存関係は存在しないものとして設計する」ルールに従い、
依存関係があるモジュール同士は1つに統合する。
ChassisModuleとDriveMotorModuleは明確に依存関係があるため、
ChassisModuleがDriveMotorModuleを内部部品として包含している。
\end{pointbox}

\begin{notebox}[Config構造体のネスト]
\texttt{ChassisConfig.motors[4]}は\texttt{DriveMotorConfig}の配列。
Config構造体は別のConfig構造体をメンバとして持てる。
これにより\texttt{ProjectConfig.h}で4輪分のピンアサインを一括定義できる。
\end{notebox}

% BleModule
\section{【応用】BleModule — スレッドセーフティパターン}

\begin{patternbox}{BleModule で学べるパターン}
\begin{itemize}
    \item volatileフラグ方式による安全なデータ受け渡し
    \item BLEコールバック（別タスク）→内部バッファ→update()→SystemData
    \item deinit()によるBLEリソースの解放
    \item 双方向通信（RX: コールバック受信、TX: update()で送信）
\end{itemize}
\end{patternbox}

\subsection{内部バッファとvolatileフラグ}

\begin{lstlisting}[style=cppstyle, caption={BleModule.h — 内部バッファ}]
class BleModule : public IModule {
private:
    // 内部バッファ（BLEコールバックスレッドが書き込む）
    uint8_t       _rxBuffer[BLE_RX_BUFFER_SIZE] = {};
    uint8_t       _rxBufferLength = 0;
    volatile bool _dataReceived   = false;  // volatileフラグ
    volatile bool _clientConnected = false;
public:
    // BLEコールバックから呼ばれる（別タスクで実行される）
    void onReceive(const uint8_t* payload, uint8_t length);
    // ...
};
\end{lstlisting}

\begin{cautionbox}[なぜvolatileが必要か]
BLEコールバックはArduinoのメインループとは別のFreeRTOSタスクで実行される。
\texttt{volatile}がないと、コンパイラが\texttt{\_dataReceived}の値をレジスタにキャッシュし、
メインループ側で変更を検出できない場合がある。
\texttt{volatile}は「この変数は外部から変更されるかもしれない」とコンパイラに伝えるキーワード。
\end{cautionbox}

\subsection{コールバック→内部バッファ→SystemData}

\begin{lstlisting}[style=cppstyle, caption={BleModule.cpp — データ受け渡しの流れ}]
// BLEコールバック（別タスクから呼ばれる）
void BleModule::onReceive(const uint8_t* payload, uint8_t length) {
    memcpy(_rxBuffer, payload, length);  // (1) バッファに書き込み
    _rxBufferLength = length;
    _dataReceived = true;                // (2) フラグを立てる
}

// update()は3フェーズループ内で呼ばれる（メインタスク）
void BleModule::update(SystemData& data) {
    data.ble.connected = _clientConnected;

    if (_dataReceived) {                             // (3) フラグ確認
        memcpy(data.ble.rxData, _rxBuffer, _rxBufferLength); // (4) コピー
        data.ble.rxLength  = _rxBufferLength;
        data.ble.newDataIn = true;
        _dataReceived = false;                       // (5) フラグクリア
    }

    // 送信リクエスト処理
    if (data.ble.sendRequest && _clientConnected && _txChar) {
        _txChar->setValue(data.ble.txData, data.ble.txLength);
        _txChar->notify();
        data.ble.sendRequest = false;
    }
}
\end{lstlisting}

\begin{pointbox}[volatileフラグ方式の安全性]
\texttt{bool}型の読み書きは原子的（1命令で完了する）。
「バッファに全データを書き込んだ後にフラグを立てる」→「フラグを確認してからバッファを読む」
の順序により、バッファの途中状態を読むことがない。
mutexやセマフォは不要で、割り込みコンテキストでも安全に使える。
\end{pointbox}

\subsection{BLEコールバックの登録}

\begin{lstlisting}[style=cppstyle, caption={BleModule.cpp — コールバッククラス}]
// BLEライブラリのコールバッククラスを継承
class BleRxCallbacks : public BLECharacteristicCallbacks {
private:
    BleModule* _module;
public:
    BleRxCallbacks(BleModule* module) : _module(module) {}
    void onWrite(BLECharacteristic* characteristic) override {
        std::string value = characteristic->getValue();
        uint8_t len = min((int)value.length(), BLE_RX_BUFFER_SIZE);
        _module->onReceive(
            reinterpret_cast<const uint8_t*>(value.data()), len);
    }
};
\end{lstlisting}

\begin{pointbox}[コールバックからモジュールへの橋渡し]
BLEライブラリはコールバッククラスの継承を要求するため、
内部クラスでラップして\texttt{BleModule::onReceive()}に橋渡しする。
コールバッククラスはモジュールのポインタを保持し、受信データをモジュールの内部バッファに渡す。
\end{pointbox}

\subsection{送信の流れ（ロジックフェーズ→update()）}

\begin{lstlisting}[style=cppstyle, caption={ロジックフェーズでの送信リクエスト}]
void applyPattern(SystemData& data) {
    // BLE受信データの処理
    if (data.ble.newDataIn) {
        // 受信データを処理...
        data.ble.newDataIn = false;  // 処理済みフラグクリア
    }

    // BLE送信（応答を返す例）
    if (needToRespond) {
        memcpy(data.ble.txData, response, responseLen);
        data.ble.txLength = responseLen;
        data.ble.sendRequest = true;  // 次のupdate()で送信される
    }
}
\end{lstlisting}

\begin{notebox}[双方向通信のデータフロー]
\textbf{受信}: BLEコールバック → 内部バッファ → update() → SystemData.ble.rxData → ロジックフェーズ\\
\textbf{送信}: ロジックフェーズ → SystemData.ble.txData → update() → BLE Notify\\
受信は入力フェーズの流れ、送信は出力フェーズの流れに対応する。
BleModuleは入力配列に登録するが、送信処理も\texttt{update()}内で行うため実質的に入出力兼用。
\end{notebox}

% WifiModule
\section{【応用】WifiModule — ステートマシンパターン}

\begin{patternbox}{WifiModule で学べるパターン}
\begin{itemize}
    \item enum classによる状態定義とswitch文による状態遷移
    \item ModuleTimerによるタイムアウト・再接続間隔の管理
    \item リクエストフラグによる外部からの制御（接続/切断指示）
    \item deinit()によるWiFiリソースの解放
\end{itemize}
\end{patternbox}

\subsection{状態定義}

\begin{lstlisting}[style=cppstyle, caption={WifiModule.h — 状態定義}]
enum class WifiState : uint8_t {
    DISCONNECTED,  // 未接続（接続リクエスト待ち）
    CONNECTING,    // 接続試行中
    CONNECTED,     // 接続済み
    RECONNECTING,  // 再接続試行中
    FAILED,        // 接続失敗（リトライ上限超過）
};
\end{lstlisting}

\begin{pointbox}[Data構造体にステートを公開]
\texttt{WifiData.state}としてenum値をSystemDataに公開することで、
ロジックフェーズが現在の接続状態に応じた判断を行える。
\texttt{isConnected}は簡易参照用のboolで、\texttt{state == CONNECTED}の省略形。
\end{pointbox}

\subsection{ステートマシンの実装}

\begin{lstlisting}[style=cppstyle, caption={WifiModule.cpp — updateInput()のステートマシン}]
void WifiModule::updateInput(SystemData& data) {
    // ステートマシン（WiFi状態の監視・自動再接続）
    switch (_state) {
        case WifiState::CONNECTING:
        case WifiState::RECONNECTING:
            if (WiFi.status() == WL_CONNECTED) {
                _state = WifiState::CONNECTED;
                _retryCount = 0;
            } else if (_stateTimer.getNowTime()
                       >= _config.connectTimeoutMs) {
                WiFi.disconnect();
                _retryCount++;
                if (_config.maxRetries > 0
                    && _retryCount >= _config.maxRetries) {
                    _state = WifiState::FAILED;
                } else {
                    _state = WifiState::RECONNECTING;
                    _stateTimer.setTime();
                }
            }
            break;

        case WifiState::CONNECTED:
            if (WiFi.status() != WL_CONNECTED) {
                _state = WifiState::RECONNECTING;
                _stateTimer.setTime();
            }
            break;
        // ...
    }

    // SystemDataに状態を反映
    data.wifi.state       = _state;
    data.wifi.isConnected = (_state == WifiState::CONNECTED);
    data.wifi.retryCount  = _retryCount;
}
\end{lstlisting}

\begin{lstlisting}[style=cppstyle, caption={WifiModule.cpp — updateOutput()のリクエスト処理}]
void WifiModule::updateOutput(SystemData& data) {
    // ロジックフェーズからのリクエスト処理
    if (data.wifi.requestConnect) {
        data.wifi.requestConnect = false;
        if (_state == WifiState::DISCONNECTED
            || _state == WifiState::FAILED) {
            _retryCount = 0;
            _startConnect();
        }
    }
    if (data.wifi.requestDisconnect) {
        data.wifi.requestDisconnect = false;
        _disconnect();
    }
}
\end{lstlisting}

\begin{pointbox}[ノンブロッキングなステートマシン]
WiFi接続は通常数秒かかるが、\texttt{delay()}で待たない。
\texttt{CONNECTING}状態では毎ループ\texttt{WiFi.status()}を確認し、
接続完了かタイムアウトかを判定する。
これにより、WiFi接続中も他のモジュール（LED、センサー等）は正常に動作し続ける。
\end{pointbox}

\begin{pointbox}[リクエストフラグパターン]
\texttt{requestConnect}はロジックフェーズからの「接続して」という指示。
BuzzerModuleの\texttt{requestPlay}と同じパターンで、
\texttt{updateOutput()}が処理後に\texttt{false}に戻す。
\end{pointbox}

\subsection{ロジックフェーズでの使用例}

\begin{lstlisting}[style=cppstyle, caption={ロジックフェーズでのWiFi制御}]
void applyPattern(SystemData& data) {
    // ボタンが押されたらWiFi接続開始
    if (data.button.justPressed
        && !data.wifi.isConnected) {
        data.wifi.requestConnect = true;
    }

    // WiFi接続状態に応じた表示
    switch (data.wifi.state) {
        case WifiState::CONNECTING:
            snprintf(data.display.line1, 48, "WiFi: Connecting...");
            break;
        case WifiState::CONNECTED:
            snprintf(data.display.line1, 48,
                     "WiFi: OK (RSSI=%d)", data.wifi.rssi);
            break;
        case WifiState::FAILED:
            snprintf(data.display.line1, 48,
                     "WiFi: Failed (%d retries)",
                     data.wifi.retryCount);
            break;
        default: break;
    }
}
\end{lstlisting}

% ─── API仕様 ───────────────────────────────
\subsection{API仕様}


\begin{apibox}{WifiModule --- WiFi接続管理モジュール}
\textbf{定義ファイル:} \texttt{lib/WifiModule/WifiModule.h / .cpp}\\
\textbf{分類:} 入出力モジュール（\texttt{updateInput()} + \texttt{updateOutput()}）\\
\textbf{対象デバイス:} ESP32 WiFi

WiFi接続のライフサイクル（接続・切断・再接続・タイムアウト）をステートマシンで自動管理するモジュール。
ロジックフェーズからリクエストフラグで接続・切断を制御する。
\end{apibox}

\subsubsection{機能概要}

\begin{itemize}
    \item 5状態のステートマシンによる接続ライフサイクル管理
    \item 自動再接続（リトライ回数制限付き）
    \item 接続タイムアウト検出
    \item RSSI（信号強度）とローカルIPの取得
    \item リクエストフラグによる接続・切断制御
    \item \texttt{deinit()}によるWiFi切断
\end{itemize}

% --- 列挙型 ---
\subsubsection{列挙型}

\begin{funcblock}{WifiState}
WiFi接続のステートマシン状態を定義する。

\begin{longtable}{l l p{6.0cm}}
\toprule
\textbf{値} & \textbf{名前} & \textbf{説明} \\
\midrule
\endhead
0 & \texttt{DISCONNECTED} & 未接続。\texttt{requestConnect}で接続開始 \\
1 & \texttt{CONNECTING}   & 接続試行中。タイムアウトまで待機 \\
2 & \texttt{CONNECTED}    & 接続済み。RSSI・IPを定期更新 \\
3 & \texttt{RECONNECTING} & 再接続試行中。切断検出後に自動遷移 \\
4 & \texttt{FAILED}       & 接続失敗。リトライ上限超過またはタイムアウト \\
\bottomrule
\end{longtable}
\end{funcblock}

% --- Config ---
\subsubsection{Config構造体}

\begin{funcblock}{WifiConfig}
WiFi接続のパラメータを定義する。

\begin{longtable}{l l l p{4.0cm}}
\toprule
\textbf{メンバ} & \textbf{型} & \textbf{制約・既定値} & \textbf{説明} \\
\midrule
\endhead
\texttt{ssid}                & \texttt{const char*}   & ---          & WiFi SSID \\
\texttt{password}            & \texttt{const char*}   & ---          & WiFiパスワード \\
\texttt{connectTimeoutMs}    & \texttt{unsigned long} & 例: 10000    & 接続タイムアウト [ms] \\
\texttt{reconnectIntervalMs} & \texttt{unsigned long} & 例: 5000     & 再接続までの待機時間 [ms] \\
\texttt{maxRetries}          & \texttt{uint8\_t}      & 例: 5（0=無制限）& 最大リトライ回数 \\
\bottomrule
\end{longtable}
\end{funcblock}

% --- Data ---
\subsubsection{Data構造体}

\begin{funcblock}{WifiData}
WiFi接続の状態と制御フラグを保持する。

\begin{longtable}{l l l p{3.8cm}}
\toprule
\textbf{メンバ} & \textbf{型} & \textbf{初期値} & \textbf{説明} \\
\midrule
\endfirsthead
\toprule
\textbf{メンバ} & \textbf{型} & \textbf{初期値} & \textbf{説明} \\
\midrule
\endhead
\midrule
\multicolumn{4}{r}{\small 次ページに続く} \\
\endfoot
\bottomrule
\endlastfoot
\multicolumn{4}{l}{\textbf{--- 状態（入力フェーズで更新） ---}} \\
\texttt{state}       & \texttt{WifiState} & \texttt{DISCONNECTED} & 現在のステートマシン状態 \\
\texttt{isConnected} & \texttt{bool}      & \texttt{false}        & 接続中フラグ（簡易参照用） \\
\texttt{retryCount}  & \texttt{uint8\_t}  & \texttt{0}            & 現在のリトライ回数 \\
\texttt{rssi}        & \texttt{int8\_t}   & \texttt{0}            & 信号強度 [dBm]（接続中のみ有効） \\
\texttt{localIp}     & \texttt{uint32\_t} & \texttt{0}            & ローカルIPアドレス（uint32\_t形式） \\
\midrule
\multicolumn{4}{l}{\textbf{--- 制御（ロジックフェーズで設定） ---}} \\
\texttt{requestConnect}    & \texttt{bool} & \texttt{false} & 接続リクエスト。\texttt{true}で接続開始 \\
\texttt{requestDisconnect} & \texttt{bool} & \texttt{false} & 切断リクエスト。\texttt{true}で切断 \\
\end{longtable}
\end{funcblock}

% --- メソッド ---
\subsubsection{メソッド}

\begin{funcblock}{WifiModule(const WifiConfig\& config)}
\begin{tabularx}{\textwidth}{l X}
\toprule
\textbf{項目} & \textbf{内容} \\
\midrule
引数   & \texttt{config} --- \texttt{WifiConfig}へのconst参照 \\
動作   & Configを\texttt{\_config}にコピー保持する \\
\bottomrule
\end{tabularx}
\end{funcblock}

\begin{funcblock}{bool init()}
\begin{tabularx}{\textwidth}{l X}
\toprule
\textbf{項目} & \textbf{内容} \\
\midrule
戻り値   & \texttt{bool} --- 常に\texttt{true} \\
動作     & WiFiモードをSTAに設定し、自動再接続を無効化する（ステートマシンで管理するため）。状態を\texttt{DISCONNECTED}に初期化 \\
\bottomrule
\end{tabularx}
\end{funcblock}

\begin{funcblock}{void updateInput(SystemData\& data)}
\begin{tabularx}{\textwidth}{l X}
\toprule
\textbf{項目} & \textbf{内容} \\
\midrule
書き込み先     & \texttt{data.wifi} \\
タイミング制御 & 毎ループ実行 \\
動作           & WiFi接続状態を確認し、ステートマシンの状態遷移を処理。\texttt{state}・\texttt{isConnected}・\texttt{rssi}・\texttt{localIp}・\texttt{retryCount}を更新 \\
\bottomrule
\end{tabularx}
\end{funcblock}

\begin{funcblock}{void updateOutput(SystemData\& data)}
\begin{tabularx}{\textwidth}{l X}
\toprule
\textbf{項目} & \textbf{内容} \\
\midrule
読み取り元     & \texttt{data.wifi} \\
タイミング制御 & 毎ループ実行 \\
動作           & \texttt{requestConnect}が\texttt{true}なら接続開始、\texttt{requestDisconnect}が\texttt{true}なら切断。各フラグを処理後に\texttt{false}にクリア \\
\bottomrule
\end{tabularx}
\end{funcblock}

\begin{funcblock}{void deinit()}
\begin{tabularx}{\textwidth}{l X}
\toprule
\textbf{項目} & \textbf{内容} \\
\midrule
動作   & WiFi接続を切断し、WiFiモードをOFFにする \\
\bottomrule
\end{tabularx}
\end{funcblock}

\begin{notebox}[ステートマシン遷移]
状態遷移の概要:

\begin{itemize}
    \item \texttt{DISCONNECTED} $\rightarrow$ \texttt{CONNECTING}（\texttt{requestConnect}）
    \item \texttt{CONNECTING} $\rightarrow$ \texttt{CONNECTED}（接続成功）
    \item \texttt{CONNECTING} $\rightarrow$ \texttt{FAILED}（タイムアウト）
    \item \texttt{CONNECTED} $\rightarrow$ \texttt{RECONNECTING}（切断検出）
    \item \texttt{CONNECTED} $\rightarrow$ \texttt{DISCONNECTED}（\texttt{requestDisconnect}）
    \item \texttt{RECONNECTING} $\rightarrow$ \texttt{CONNECTED}（接続成功）
    \item \texttt{RECONNECTING} $\rightarrow$ \texttt{FAILED}（リトライ超過）
    \item \texttt{FAILED} $\rightarrow$ \texttt{CONNECTING}（\texttt{requestConnect}、リトライカウントリセット）
\end{itemize}
\end{notebox}

% --- 使用例 ---
\subsubsection{使用例}

\begin{lstlisting}[style=cppstyle, caption={WifiModuleの使用例}]
// ProjectConfig.h
const WifiConfig WIFI_CONFIG = {
    .ssid                = "YourSSID",
    .password            = "YourPassword",
    .connectTimeoutMs    = 10000,  // 10秒タイムアウト
    .reconnectIntervalMs = 5000,   // 5秒後に再接続
    .maxRetries          = 5,      // 最大5回リトライ
};

// ロジックフェーズでの使用
void applyPattern(SystemData& data) {
    // ボタンでWiFi接続トグル
    if (data.button.justPressed) {
        if (data.wifi.isConnected) {
            data.wifi.requestDisconnect = true;
        } else {
            data.wifi.requestConnect = true;
        }
    }
    // 接続中はRSSIを表示
    if (data.wifi.isConnected) {
        snprintf(data.display.line3, 48,
                 "RSSI: %d dBm", data.wifi.rssi);
    }
}
\end{lstlisting}

% EncoderModule
\section{【応用】EncoderModule — 割り込み連携パターン}

\begin{patternbox}{EncoderModule で学べるパターン}
\begin{itemize}
    \item ハードウェア割り込み（ISR）とvolatile変数の連携
    \item IRAM\_ATTR による割り込みハンドラのRAM配置
    \item ISR→volatile変数→updateInput()→SystemDataのデータフロー
    \item deinit()による割り込みの解除
\end{itemize}
\end{patternbox}

\subsection{割り込みハンドラの設計}

\begin{lstlisting}[style=cppstyle, caption={EncoderModule.h — 割り込み関連}]
class EncoderModule : public IModule {
private:
    EncoderConfig _config;
    ModuleTimer   _sampleTimer;

    // ISRが更新する変数（volatile必須）
    volatile int32_t _isrCount = 0;

    // updateInput()内で使う前回値
    int32_t _prevCount = 0;

    // ISRはstatic関数でなければattachInterruptに渡せない
    static void IRAM_ATTR _isrHandlerA();
    static EncoderModule* _instance;

    void IRAM_ATTR _onPulseA();
    // ...
};
\end{lstlisting}

\begin{notebox}[なぜstaticが必要なのか？ — 初学者向け解説]
C++のクラスのメンバ関数（例: \texttt{\_onPulseA()}）は、内部的に「どのオブジェクトの関数か」という情報（thisポインタ）を持っています。しかし、\texttt{attachInterrupt()}が受け付けるのはC言語形式の「普通の関数ポインタ」だけで、thisポインタを渡す仕組みがありません。

そこで以下の2段構えで解決します。

\begin{enumerate}
    \item \textbf{static関数}を用意する — static関数はthisポインタを必要としないので、割り込みに登録できる
    \item \textbf{staticなインスタンスポインタ}を用意する — static関数の中から「どのオブジェクトか」を知るための橋渡し
\end{enumerate}

\begin{lstlisting}
// 1. staticポインタ: 「どのエンコーダか」を覚えておく
static EncoderModule* _instance;

// 2. static関数: 割り込みに登録できる
static void _isrHandlerA() {
    _instance->_onPulseA();  // ここで実際の処理に橋渡し
}

// init()で登録
_instance = this;
attachInterrupt(pin, _isrHandlerA, RISING);
\end{lstlisting}

これは「C++のオブジェクト指向とC言語の割り込みの橋渡し」パターンです。割り込みを使わない通常のモジュールでは不要なので、初めて見る場合はこのパターンだけ覚えれば十分です。
\end{notebox}

\begin{cautionbox}[ISR（割り込みサービスルーチン）の制約]
\begin{itemize}
    \item \texttt{attachInterrupt()}に渡す関数は\textbf{static関数}でなければならない
    \item ISR内では\texttt{Serial.print()}やメモリ確保等の重い処理は禁止
    \item ESP32等では\texttt{IRAM\_ATTR}を付けてRAMに配置する必要がある（ボード依存）
    \item ISRから参照する変数は\texttt{volatile}を付ける
\end{itemize}
\end{cautionbox}

\subsection{ISR→volatile→updateInput()の流れ}

\begin{lstlisting}[style=cppstyle, caption={EncoderModule.cpp — ISRとupdate()}]
// ISRハンドラ（IRAM_ATTR = RAM配置、高速アクセス）
void IRAM_ATTR EncoderModule::_isrHandlerA() {
    if (_instance) {
        _instance->_onPulseA();
    }
}

void IRAM_ATTR EncoderModule::_onPulseA() {
    // A相の立ち上がり時にB相を読み取って方向を判定
    if (digitalRead(_config.pinB) == HIGH) {
        _isrCount++;   // CW
    } else {
        _isrCount--;   // CCW
    }
}

// updateInput()は入力フェーズで呼ばれる
void EncoderModule::updateInput(SystemData& data) {
    // volatile変数を一度だけ読み取り
    int32_t currentCount = _isrCount;

    data.encoder.count = currentCount;

    // 速度計算（一定間隔で実行）
    if (_sampleTimer.getNowTime() >= _config.sampleIntervalMs) {
        unsigned long elapsed = _sampleTimer.getNowTime();
        _sampleTimer.setTime();

        int32_t delta = currentCount - _prevCount;
        _prevCount = currentCount;

        float revolutions = (float)delta / _config.pulsesPerRev;
        float seconds = (float)elapsed / 1000.0f;
        data.encoder.rpm = (seconds > 0)
            ? (revolutions / seconds * 60.0f) : 0.0f;
        data.encoder.isValid = true;
    }
}
\end{lstlisting}

\begin{pointbox}[volatile変数を一度だけ読み取る]
\texttt{int32\_t currentCount = \_isrCount;}で値をローカル変数にコピーしてから使う。
volatile変数を関数内で複数回参照すると、参照のたびに値が変わる可能性がある。
一度だけ読み取ることで、同一updateInput()内での値の一貫性を保証する。
\end{pointbox}

\begin{pointbox}[BleModuleとの共通パターン]
BleModule（コールバック）とEncoderModule（ISR）は、
「メインタスク外からのデータ受け渡し」という同じ問題を解決している。
どちらも\textbf{volatile変数（またはフラグ）}を介して安全にデータを渡す。
BLEはバッファ+フラグ、エンコーダはカウンタ変数という違いはあるが、原理は同じ。
\end{pointbox}

\subsection{deinit()で割り込みを解除}

\begin{lstlisting}[style=cppstyle, caption={EncoderModule.cpp — deinit()}]
void EncoderModule::deinit() {
    detachInterrupt(digitalPinToInterrupt(_config.pinA));
    _instance = nullptr;
    Serial.println("[Encoder] deinit OK");
}
\end{lstlisting}

\begin{pointbox}[割り込み解除の重要性]
割り込みを解除しないままスリープに入ると、
スリープ中にも割り込みが発生してウェイクアップの原因になる可能性がある。
また、\texttt{\_instance = nullptr}にすることで、
deinit後にISRが呼ばれてもnullチェックで安全にスキップされる。
\end{pointbox}

\begin{notebox}[複数インスタンスへの拡張]
現在の実装はstaticポインタ1つ（\texttt{\_instance}）のため、
エンコーダモジュールは1つしか使えない。
複数使う場合は\texttt{static EncoderModule* \_instances[MAX]}の配列に拡張し、
チャネルIDでインデックスする方式に変更する。
\end{notebox}

% ─── API仕様 ───────────────────────────────
\subsection{API仕様}


\begin{apibox}{EncoderModule --- ロータリーエンコーダ入力モジュール}
\textbf{定義ファイル:} \texttt{lib/EncoderModule/EncoderModule.h / .cpp}\\
\textbf{分類:} 入力モジュール（\texttt{updateInput()}のみオーバーライド）\\
\textbf{対象デバイス:} インクリメンタル型ロータリーエンコーダ（2相出力）

ハードウェア割り込みでA相パルスをカウントし、B相の状態から回転方向を判定する。
周期的にRPM（回転速度）を計算して\texttt{SystemData}に反映する。
\end{apibox}

\subsubsection{機能概要}

\begin{itemize}
    \item ハードウェア割り込み（A相）によるパルスカウント
    \item B相によるCW/CCW方向判定
    \item 周期的RPM計算（\texttt{ModuleTimer}ベース）
    \item 累積パルス数の保持（正:CW、負:CCW）
    \item \texttt{deinit()}による割り込み解除
\end{itemize}

% --- Config ---
\subsubsection{Config構造体}

\begin{funcblock}{EncoderConfig}
エンコーダの接続ピンと計測パラメータを定義する。

\begin{longtable}{l l l p{4.8cm}}
\toprule
\textbf{メンバ} & \textbf{型} & \textbf{制約・既定値} & \textbf{説明} \\
\midrule
\endhead
\texttt{pinA}             & \texttt{uint8\_t}       & ---       & エンコーダA相ピン（割り込み対応ピンを使用すること） \\
\texttt{pinB}             & \texttt{uint8\_t}       & ---       & エンコーダB相ピン \\
\texttt{pulsesPerRev}     & \texttt{uint16\_t}      & 例: 20    & 1回転あたりのパルス数 \\
\texttt{sampleIntervalMs} & \texttt{unsigned long}  & 例: 50    & RPM計算間隔 [ms] \\
\bottomrule
\end{longtable}
\end{funcblock}

% --- Data ---
\subsubsection{Data構造体}

\begin{funcblock}{EncoderData}
エンコーダの計測結果を保持する。

\begin{longtable}{l l l p{5.0cm}}
\toprule
\textbf{メンバ} & \textbf{型} & \textbf{初期値} & \textbf{説明} \\
\midrule
\endhead
\texttt{count}   & \texttt{int32\_t} & \texttt{0}     & 累積パルス数。正:CW（時計回り）、負:CCW（反時計回り） \\
\texttt{rpm}     & \texttt{float}    & \texttt{0.0f}  & 回転速度 [rpm] \\
\texttt{isValid} & \texttt{bool}     & \texttt{false} & RPM計算が1回以上行われたら\texttt{true} \\
\bottomrule
\end{longtable}
\end{funcblock}

% --- メソッド ---
\subsubsection{メソッド}

\begin{funcblock}{EncoderModule(const EncoderConfig\& config)}
\begin{tabularx}{\textwidth}{l X}
\toprule
\textbf{項目} & \textbf{内容} \\
\midrule
引数   & \texttt{config} --- \texttt{EncoderConfig}へのconst参照 \\
動作   & Configを\texttt{\_config}にコピー保持する \\
\bottomrule
\end{tabularx}
\end{funcblock}

\begin{funcblock}{bool init()}
\begin{tabularx}{\textwidth}{l X}
\toprule
\textbf{項目} & \textbf{内容} \\
\midrule
戻り値   & \texttt{bool} --- 常に\texttt{true} \\
動作     & A相・B相ピンを\texttt{INPUT\_PULLUP}で設定。A相にRISINGエッジの割り込みハンドラを登録。シングルインスタンス参照（\texttt{\_instance}）を自身に設定 \\
\bottomrule
\end{tabularx}
\end{funcblock}

\begin{funcblock}{void updateInput(SystemData\& data)}
\begin{tabularx}{\textwidth}{l X}
\toprule
\textbf{項目} & \textbf{内容} \\
\midrule
書き込み先     & \texttt{data.encoder} \\
タイミング制御 & \texttt{ModuleTimer}で\texttt{\_config.sampleIntervalMs}ごとにRPMを計算 \\
動作           & ISRが更新した\texttt{volatile}カウンタを読み取り、前回値との差分からRPMを計算。\texttt{count}・\texttt{rpm}・\texttt{isValid}を更新 \\
\bottomrule
\end{tabularx}
\end{funcblock}

\begin{funcblock}{void deinit()}
\begin{tabularx}{\textwidth}{l X}
\toprule
\textbf{項目} & \textbf{内容} \\
\midrule
動作   & A相ピンの割り込みを\texttt{detachInterrupt()}で解除する \\
\bottomrule
\end{tabularx}
\end{funcblock}

\begin{notebox}[RPM計算式]
RPM計算は以下の式で行う:

\vspace{0.3em}
$$RPM = \frac{\Delta count}{pulsesPerRev} \times \frac{60000}{sampleIntervalMs}$$
\vspace{0.3em}

\noindent ここで $\Delta count$ はサンプリング間隔内のパルス数変化量である。
\end{notebox}

\begin{cautionbox}[シングルインスタンス制約]
\texttt{attachInterrupt()}にstatic関数を渡すため、内部で\texttt{static EncoderModule* \_instance}を使用している。
同時に使用できるEncoderModuleのインスタンスは\textbf{1つのみ}である。
複数エンコーダが必要な場合はISRディスパッチの仕組みを拡張する必要がある。
\end{cautionbox}

% --- 使用例 ---
\subsubsection{使用例}

\begin{lstlisting}[style=cppstyle, caption={EncoderModuleの使用例}]
// ProjectConfig.h
const EncoderConfig ENCODER_CONFIG = {
    .pinA             = 39,  // GPIO39: A相
    .pinB             = 40,  // GPIO40: B相
    .pulsesPerRev     = 20,  // 1回転20パルス
    .sampleIntervalMs = 50,  // 50msごとに速度計算
};

// ロジックフェーズでの使用
void applyPattern(SystemData& data) {
    if (data.encoder.isValid) {
        float speed = data.encoder.rpm;
        // 回転速度に応じた制御
    }
}
\end{lstlisting}


\end{document}
